% Introduction
\section{Introduction}\label{sec:intro}

Hit-to-lead (H2L) is structurally and operationally distinct from de novo drug discovery. Medicinal chemists begin from a measured hit with a known binding mode and traverse a small number of structure-activity-relationship (SAR) steps toward a candidate that is both potent and developable. Each generated molecule is judged on many axes (predicted affinity, drug-likeness, synthetic tractability, selectivity, metabolic stability, physicochemical fit to the binding-site environment, and absence of toxicophores), only a small fraction of which can be encoded as scalar reward terms in a generative loop. For \emph{covalent} programs an additional, hard constraint applies: the candidate must retain a defined electrophilic warhead positioned for the canonical thia-Michael transition state \cite{lonsdale2018structure,burgi1974stereochemistry}. This constraint is uniquely unforgiving: a molecule that loses its warhead is no longer a covalent inhibitor at all, and a warhead presented in the wrong orientation may either fail to reach the intended cysteine or react promiscuously elsewhere on the proteome. The covalent generative problem is therefore not merely ``generate a potent ligand'' but ``generate a potent ligand whose specific electrophilic atom can reach the target nucleophile in the productive transition-state geometry''.

The covalent-aware design problem decomposes into a short stack of progressively harder constraints (Table~\ref{tab:covalent_levels}). \textbf{P1} just requires a valid drug-like molecule, which any pretrained chemical language model solves. \textbf{P2} adds the requirement that the molecule contains the warhead substructure; methods solve it either by training-time enrichment (covFT) or by hard-freezing the warhead atoms at the SMILES level (LibInvent and similar scaffold-decoration models). \textbf{P3} additionally requires that the warhead adopts the planar enone geometry needed for $\pi$-conjugation, which is the prerequisite for the alkene to be electrophilic at all; template-freezing reaches P3 by extending the frozen fragment to include the atoms that lock the dihedral, but this comes at the cost of generation flexibility (more atoms frozen $\Rightarrow$ less novel chemistry). \textbf{P4} is the only level that requires pocket-specific information: the warhead's $\beta$-carbon must physically reach the target cysteine in the productive transition-state orientation. A planar warhead in the wrong corner of the pocket is electronically productive but physically unreachable, so P3 does not imply P4.

\begin{table}[!t]
\centering
\caption{Difficulty levels for covalent-aware generation. Each row adds a constraint to the row above; the rightmost column indicates the minimum input information required to satisfy it.}
\label{tab:covalent_levels}
\renewcommand{\arraystretch}{1.5}
\setlength{\tabcolsep}{8pt}
\small
\begin{tabular}{|c|p{0.30\linewidth}|p{0.26\linewidth}|p{0.22\linewidth}|}
\hline
\textbf{Level} & \textbf{Constraint} & \textbf{Methods that achieve it} & \textbf{Required signal} \\
\hline
P1 & Valid drug-like ligand.                                                          & mol2mol prior \cite{loeffler2024reinvent4}.                                & ligand sequence. \\
\hline
P2 & + warhead substructure retained.                                                 & covFT; LibInvent \cite{fialkova2022libinvent} with frozen warhead atoms.   & ligand sequence. \\
\hline
P3 & + warhead in planar productive conformation (s-cis or s-trans, $\pi$-conjugated). & LibInvent with frozen scaffold-plus-warhead, at the cost of generation flexibility. & ligand sequence + 3D conformer. \\
\hline
P4 & + warhead reaches the target nucleophile in the pocket.                          & M1a (this work).                                                           & ligand sequence + pocket + 3D pose. \\
\hline
\end{tabular}
\end{table}

Existing generative-chemistry tools address subsets of this stack but do not jointly enforce all four. SMILES sequence generators, including the mol2mol prior used in this work \cite{loeffler2024reinvent4}, the original REINFORCE-style chemical RL of Olivecrona et al.\ \cite{olivecrona2017reinvent}, and substructure-constrained de novo generators \cite{blaschke2020reinvent2}, reach P1--P2 and can be steered with substructure filters but are pose-agnostic. Scaffold-decoration tools such as LibInvent \cite{fialkova2022libinvent} reach P3 by template-freezing the warhead atoms, but at the cost of restricting generation to a fixed parent scaffold. Pocket-aware 3D-generative methods such as DiffSBDD \cite{schneuing2024diffsbdd}, Pocket2Mol \cite{peng2022pocket2mol}, and TargetDiff \cite{guan2023targetdiff} produce ligand atoms and coordinates inside the binding cavity (analogous to a non-covalent P4) but face a separate inference-time difficulty: converting the predicted atom positions back to a SMILES requires post-hoc bond perception from coordinates, which is heuristic and lossy, and frequently produces invalid or chemically implausible molecules. None of these methods rewards an electrophilic warhead, let alone its productive orientation. Structure-prediction cofold methods \cite{passaro2025boltz2} can validate poses at P4 quality but are too slow to drive an online generative policy. No single existing tool starts from a measured hit and jointly achieves P2--P4 while generating the full molecular scaffold.

\covagen{} bridges this gap. It is a hit-to-lead framework built on a \emph{3D-covalent-aware architecture and fine-tuning} of a SMILES language model: a covalent-adapted SMILES decoder is extended with a small pocket and warhead-pose conditioning module so that the decoder can generate molecules whose electrophilic atom is presented in the geometry the target nucleophile requires. The policy is driven by a multi-objective RL loop that jointly rewards predicted potency, warhead retention, and drug-likeness; we additionally explore preference-learning variants (DPO over Boltz cofold quality) as a complementary training signal.

The contributions of this work are as follows.
\begin{enumerate}
  \item \textbf{A covalent-aware generative model} trained and tuned to allow generation conditioned on warhead--pocket geometry, including a pocket and warhead-pose-conditioned decoder (M1a) and preference-learning variants derived from Boltz cofold quality.
  \item \textbf{A hit-to-lead framework} (\covagen{}) combining the covalent-aware prior with multi-objective reinforcement learning that rewards predicted potency, warhead retention, and drug-likeness in concert.
  \item Prospective wet-lab validation of selected candidates is in flight at a partner site; the ultimate test of the pipeline is whether any of these candidates shows useful potency.
\end{enumerate}
